\documentclass[../main.tex]{subfiles}
\graphicspath{{\subfix{../img/}}{\subfix{../architectures images/}}}
\begin{document}
\chapter{Sprint 4 : Gestion des Coachs et des S\'{e}ances}

%──────────────────────────────────────────────
\section{Introduction}
%──────────────────────────────────────────────

Apr\`{e}s l'acc\`{e}s physique et la gestion des abonnements, ce sprint apporte la dimension collective de la salle : les coachs, les disciplines et les s\'{e}ances. Un g\'{e}rant peut planifier une s\'{e}ance, inviter des membres par e-mail, et les membres pointent leur pr\'{e}sence avec leur t\'{e}l\'{e}phone en scannant un code \acrshort{qr}. Trois types de s\'{e}ances couvrent les diff\'{e}rents cas --- ouvertes \`{a} tous, sur r\'{e}servation ou sur invitation --- avec un suivi de paiement int\'{e}gr\'{e} pour les coachs.

%──────────────────────────────────────────────
\section{Sprint Backlog}
%──────────────────────────────────────────────

Le tableau~\ref{tab:backlog_sprint4} pr\'{e}sente le backlog du Sprint 4 d\'{e}taill\'{e} par \'{e}pic, nature et user stories.

\begin{longtable}{|p{2.5cm}|p{2cm}|p{4cm}|p{5.5cm}|}
\hline
\textbf{Epic} & \textbf{Nature} & \textbf{User Story} & \textbf{Description} \\
\hline
\endfirsthead
\hline
\textbf{Epic} & \textbf{Nature} & \textbf{User Story} & \textbf{Description} \\
\hline
\endhead
\hline
\endfoot

\multirow{3}{*}{\parbox{2.5cm}{Gestion des coachs}} & \acrshort{crud} & En tant qu'administrateur, je veux g\'{e}rer les coachs & Cr\'{e}er, modifier, consulter et supprimer un coach avec ses informations : nom, date de naissance, profession, genre, CIN, tarif horaire, statut actif, identifiant de carte et code PIN. \\
\cline{2-4}
 & \acrshort{crud} & En tant qu'administrateur, je veux associer une image au coach & Associer une photo de profil au coach via un identifiant d'image (imageId). \\
\cline{2-4}
 & \acrshort{crud} & En tant qu'administrateur, je veux activer ou d\'{e}sactiver un coach & G\'{e}rer le statut \texttt{isEnabled} pour contr\^{o}ler la disponibilit\'{e} du coach. \\
\hline

\multirow{2}{*}{\parbox{2.5cm}{Gestion des disciplines}} & \acrshort{crud} & En tant qu'administrateur, je veux g\'{e}rer les disciplines & Cr\'{e}er et g\'{e}rer des disciplines (titre, image) associ\'{e}es \`{a} un partenaire sp\'{e}cifique. \\
\cline{2-4}
 & Consultation & En tant que membre, je veux voir les disciplines disponibles & Consulter la liste des disciplines propos\'{e}es par la salle de sport. \\
\hline

\multirow{3}{*}{\parbox{2.5cm}{Planification des s\'{e}ances}} & \acrshort{crud} & En tant qu'administrateur, je veux planifier des s\'{e}ances & Cr\'{e}er des s\'{e}ances avec date, heure de d\'{e}but, heure de fin, type (OPEN, RESERVATION, CLOSE) et nombre maximum de places pour le type RESERVATION. \\
\cline{2-4}
 & \acrshort{crud} & En tant qu'administrateur, je veux g\'{e}rer les cr\'{e}neaux des coachs & D\'{e}finir les disponibilit\'{e}s horaires des coachs pour la planification des s\'{e}ances. \\
\cline{2-4}
 & Consultation & En tant que membre, je veux consulter le planning des s\'{e}ances & Visualiser le planning hebdomadaire des s\'{e}ances avec leurs types et horaires. \\
\hline

\multirow{2}{*}{\parbox{2.5cm}{Invitations aux s\'{e}ances}} & Fonctionnalit\'{e} & En tant qu'administrateur, je veux inviter des membres \`{a} une s\'{e}ance & Envoyer des invitations \`{a} des membres sp\'{e}cifiques pour participer \`{a} une s\'{e}ance. \\
\cline{2-4}
 & Notification & En tant que membre, je veux recevoir une notification par e-mail & Recevoir un e-mail de notification lors d'une invitation \`{a} une s\'{e}ance. \\
\hline

\multirow{1}{*}{\parbox{2.5cm}{Pointage \acrshort{qr}}} & Fonctionnalit\'{e} & En tant que membre, je veux pointer ma pr\'{e}sence via un code \acrshort{qr} & Scanner un code \acrshort{qr} de s\'{e}ance pour confirmer la pr\'{e}sence lors d'une session. \\
\hline

\multirow{2}{*}{\parbox{2.5cm}{Paiements des s\'{e}ances}} & Fonctionnalit\'{e} & En tant qu'administrateur, je veux g\'{e}rer les paiements des s\'{e}ances & Enregistrer les paiements pour les s\'{e}ances de coaching payantes. \\
\cline{2-4}
 & Consultation & En tant qu'administrateur, je veux suivre les paiements des coachs & Consulter l'historique des paiements et g\'{e}n\'{e}rer les factures de r\'{e}mun\'{e}ration des coachs. \\
\hline
\caption{Sprint Backlog -- Sprint 4}
\label{tab:backlog_sprint4}
\end{longtable}

%──────────────────────────────────────────────
\section{Analyse}
%──────────────────────────────────────────────

\subsection{Diagramme de cas d'utilisation}

La figure~\ref{fig:sprint4-usecase} montre la r\'{e}partition des actions : l'administrateur pilote tout le cycle coach/s\'{e}ance, le membre consulte le planning et pointe sa pr\'{e}sence.

\begin{figure}[H]
\centering
\begin{tikzpicture}[scale=0.85, transform shape]

% System box
\node[systembox, minimum width=10cm, minimum height=14cm] (sys) at (5.5, -6) {};
\node[uctitle] at (sys.north) {Syst\`{e}me TrustGym -- Coachs et S\'{e}ances};

% Actors (all on the left)
\stickfigure{(-2.5, -3)}{Admin}
\stickfigure{(-2.5, -9)}{Membre}

% Use cases
\node[usecase] (uc1) at (5.5, 0) {G\'{e}rer les coachs};
\node[usecase] (uc2) at (5.5, -1.8) {G\'{e}rer les disciplines};
\node[usecase] (uc3) at (5.5, -3.6) {Planifier une s\'{e}ance};
\node[usecase] (uc4) at (5.5, -5.4) {G\'{e}rer les cr\'{e}neaux coachs};
\node[usecase] (uc5) at (5.5, -7.2) {Inviter des membres};
\node[usecase] (uc6) at (5.5, -9) {G\'{e}rer les paiements coachs};
\node[usecase] (uc7) at (5.5, -10.8) {Consulter le planning};
\node[usecase] (uc8) at (5.5, -12.6) {Pointer via code QR};

% Admin links
\draw (-1, -3) -- (uc1.west);
\draw (-1, -3) -- (uc2.west);
\draw (-1, -3) -- (uc3.west);
\draw (-1, -3) -- (uc4.west);
\draw (-1, -3) -- (uc5.west);
\draw (-1, -3) -- (uc6.west);

% Membre links
\draw (-1, -9) -- (uc7.west);
\draw (-1, -9) -- (uc8.west);

\end{tikzpicture}
\caption{Diagramme de cas d'utilisation -- Sprint 4}
\label{fig:sprint4-usecase}
\end{figure}

\subsection{Cas d'utilisation cibl\'{e}s}

Format d\'{e}fini en section~\ref{sec:uc_format}. Ce sprint couvre : \textbf{G\'{e}rer les coachs} (Admin), \textbf{Planifier une s\'{e}ance} (Admin --- contr\^{o}le de disponibilit\'{e} via \texttt{CoachTimeSlot}, types OPEN/RESERVATION/CLOSE avec \texttt{maxSlots}), et \textbf{Pointer via code \acrshort{qr}} (Membre --- scan via \texttt{expo-camera}, validation cumul\'{e}e session active + inscription).

%──────────────────────────────────────────────
\section{R\'{e}alisation}
%──────────────────────────────────────────────

\subsection{Backend}

\begin{longtable}{|>{\footnotesize}p{4.5cm}|>{\footnotesize}p{3.6cm}|p{5.6cm}|}
\hline
\textbf{Module} & \textbf{Entit\'{e}s cl\'{e}s} & \textbf{R\^{o}le} \\
\hline
\endfirsthead
\hline
\endfoot
\texttt{CoachModule} & \texttt{Coach}, \texttt{CoachImage} & \acrshort{crud} coach (\texttt{hourlyFee}, \texttt{isEnabled}, \texttt{cardId}, \texttt{pin}). \\
\hline
\texttt{DisciplineModule} & \texttt{Discipline} & Disciplines sportives par partenaire. \\
\hline
\texttt{CoachSessionModule} & \texttt{Session} & S\'{e}ances avec types \texttt{OPEN} (acc\`{e}s libre), \texttt{RESERVATION} (\texttt{maxSlots}) et \texttt{CLOSE} (sur invitation). \\
\hline
\texttt{CoachTimeSlotModule} & \texttt{CoachTimeSlot}, \texttt{TimeSlotHours} & Disponibilit\'{e}s hebdomadaires r\'{e}currentes ; v\'{e}rification automatique lors de la planification. \\
\hline
\texttt{SessionInvitationModule} (+\texttt{EmailModule}) & \texttt{SessionInvitation} & Invitations cibl\'{e}es par e-mail (templates Nodemailer). \\
\hline
\texttt{SessionQRCodeModule} & \texttt{SessionQRCode} & Code \acrshort{qr} unique par s\'{e}ance pour le pointage mobile. \\
\hline
\texttt{SessionPaymentModule} & \texttt{SessionPayment} & R\'{e}mun\'{e}ration des coachs et historique des paiements. \\
\hline
\caption{Modules backend du Sprint 4}
\end{longtable}

\subsection{Diagramme de s\'{e}quence}

La figure~\ref{fig:sprint4-sequence} d\'{e}taille le parcours d'un scan \acrshort{qr} : trois v\'{e}rifications c\^{o}t\'{e} API avant de confirmer la pr\'{e}sence.

\begin{figure}[H]
\centering
\begin{tikzpicture}[scale=0.8, transform shape]

% Participants
\node[sdactor] (mob) at (0, 0) {App Mobile};
\node[sdactor] (cam) at (3.5, 0) {expo-camera};
\node[sdactor] (api) at (7, 0) {NestJS API};
\node[sdactor] (db) at (11, 0) {PostgreSQL};

% Lifelines
\draw[dashed] (0, -0.35) -- (0, -12);
\draw[dashed] (3.5, -0.35) -- (3.5, -12);
\draw[dashed] (7, -0.35) -- (7, -12);
\draw[dashed] (11, -0.35) -- (11, -12);

% Messages
\draw[sdmessage] (0, -1) -- node[above, font=\scriptsize] {Ouvrir scanner} (3.5, -1);
\draw[sdreturn] (3.5, -2) -- node[above, font=\scriptsize] {QR data (sessionId)} (0, -2);

\draw[sdmessage] (0, -3.5) -- node[above, font=\scriptsize, text width=4cm, align=center] {POST /api/session-qr/verify} (7, -3.5);
\node[font=\scriptsize\itshape, text=gray] at (3.5, -4.2) {Bearer memberToken};

\draw[sdmessage] (7, -5) -- node[above, font=\scriptsize, text width=2.5cm, align=center] {V\'{e}rifier session active} (11, -5);
\draw[sdreturn] (11, -6) -- node[above, font=\scriptsize] {Session trouv\'{e}e} (7, -6);

\draw[sdmessage] (7, -7) -- node[above, font=\scriptsize, text width=2.5cm, align=center] {V\'{e}rifier inscription} (11, -7);
\draw[sdreturn] (11, -8) -- node[above, font=\scriptsize] {Membre inscrit} (7, -8);

\draw[sdmessage] (7, -9) -- node[above, font=\scriptsize, text width=2.5cm, align=center] {INSERT pr\'{e}sence} (11, -9);
\draw[sdreturn] (11, -10) -- node[above, font=\scriptsize] {OK} (7, -10);

\draw[sdreturn] (7, -11) -- node[above, font=\scriptsize] {Pr\'{e}sence confirm\'{e}e} (0, -11);

\end{tikzpicture}
\caption{Diagramme de s\'{e}quence -- V\'{e}rification de pr\'{e}sence par code QR}
\label{fig:sprint4-sequence}
\end{figure}

\subsection{Interfaces r\'{e}alis\'{e}es (web et mobile)}

Pages web : \texttt{coaches/}, \texttt{cours/} (disciplines), \texttt{seances/} (cr\'{e}ation et listing color\'{e} par type), \texttt{planning/} (calendrier), \texttt{paiement/} (r\'{e}mun\'{e}rations). C\^{o}t\'{e} mobile : \texttt{PlanningView} avec code couleur par type, scanner \acrshort{qr} via \texttt{expo-camera}.

\begin{figure}[H]
  \centering
  \includegraphics[width=\textwidth, height=0.55\textheight, keepaspectratio]{capture/planning-seances.png}
  \caption{Calendrier hebdomadaire des s\'{e}ances planifi\'{e}es}
  \label{fig:sprint4-planning}
\end{figure}

\begin{figure}[H]
  \centering
  \includegraphics[width=\textwidth, height=0.55\textheight, keepaspectratio]{capture/Scanner QR sur mobile.png}
  \caption{Scan du code QR pour le pointage de pr\'{e}sence}
  \label{fig:sprint4-qr-scan}
\end{figure}

%──────────────────────────────────────────────
\section{Conclusion}
%──────────────────────────────────────────────

TrustGym dispose maintenant d'une offre de coaching structur\'{e}e : profils de coachs, disciplines, planning hebdomadaire, invitations et pointage \acrshort{qr}. Le Sprint 5 quitte le terrain des s\'{e}ances pour celui des produits, avec la boutique et la gestion de stock.

\end{document}
